% Frozen candidate proof for
% p21p3_curve_rigidity_nonreduced_addendum_statement_20260827.tex.
% Date: 27 August 2026.
% This packet is not integrated into the live book.

\begin{proof}
The reduction has the same underlying topological space and completed
residue fields as \(X\).  Affinoid charts pass to quotients by nilpotent
ideals, so goodness, separatedness, strictness, and pure dimension one pass
to \(X_{\mathrm{red}}\).  The immersion \(\iota\) is finite and hence has
empty relative boundary by \cite[\S11.4.4]{CLD25}.  The
boundary-composition identity
\cite[(11.4.1.3)]{CLD25} therefore shows that
\(X_{\mathrm{red}}\) is boundaryless.  Thus \(X_{\mathrm{red}}\) satisfies
the ambient hypotheses of \cref{prop:p21p3-curve-rigidity}.

\smallskip
\emph{Forms, currents and psh functions.}
By \cref{prop:ch13-pullback-properties}\,
\cref{itm:ch13-pb:red}, proved from
\cite[Proposition~15.1.3(b) and the following form variant]{CLD25}, pullback
gives an isomorphism

\begin{equation}
\label{eq:p21p3-nr-form-isomorphism}
       \iota^*:A_X^{p,q}\xrightarrow{\ \sim\ }
                       A_{X_{\mathrm{red}}}^{p,q}.
\end{equation}

The point needed for currents is topological, not merely algebraic.  Recall
that a datum for the compact-test topology consists of a compact support and
a finite family of moment charts.  Restriction sends such a datum on \(X\)
to the same datum on \(X_{\mathrm{red}}\).  Conversely, the proof of the
cited almost-good result applies because \(X\) is good: at each point of the
compact support, every finite tuple of invertible functions defining a
reduced moment chart lifts to a tuple of unit germs on \(X\) having the same
absolute values on an ordinary open neighbourhood.  Shrinking these
neighbourhoods and using compactness gives a finite ordinary-open refinement
on which all moment tuples lift.  Taking product moment maps on overlaps if
necessary gives a finite lifted chart datum on \(X\).  Restriction and these
finite lifted refinements thus give cofinal systems of compact chart data for
the final locally convex topologies on the two sides of
\eqref{eq:p21p3-nr-form-isomorphism}.

A lifted invertible analytic function has the same absolute value at every
point.  After the preceding finite refinement, the tropical images,
Lagerberg representatives, polyhedral refinements, and every
\(C^r\)-seminorm are consequently the same.  Cofinality now shows that

\begin{equation}
\label{eq:p21p3-nr-test-homeomorphism}
 \iota^*:A^{p,q}_{X,c}\xrightarrow{\ \sim\ }
                         A^{p,q}_{X_{\mathrm{red}},c}
\end{equation}

is a homeomorphism.  The same chart description proves compatibility with
wedge, \(J\), the differentials, supports, and positivity of forms.  Taking
continuous duals in \eqref{eq:p21p3-nr-test-homeomorphism} gives an
isomorphism

\begin{equation}
\label{eq:p21p3-nr-current-isomorphism}
 \iota_*:D^{a,b}(X_{\mathrm{red}})
              \xrightarrow{\ \sim\ }D^{a,b}(X),
 \qquad
 \langle\iota_*S,\alpha\rangle
       =\langle S,\iota^*\alpha\rangle .
\end{equation}

It preserves positivity because the relevant test-form cones agree, and
it preserves closedness because \(\iota^*\) commutes with
\(\mathrm d'\) and \(\mathrm d''\).  Multiplication by a smooth function
commutes with \(\iota_*\) directly by transposition.  If \(S\) is positive,
the Borel product of \cite[Example~7.2.12]{CLD25} is defined through the
measures \(\mu_{S,\omega}\) attached to positive test forms \(\omega\).
Under \(\iota^*\), the positive test forms and these attached measures agree.
Thus \(\iota_*(fS)=f\iota_*S\) for every locally bounded Borel function
\(f\).  If \(S=S_1-S_2\) locally with \(S_1,S_2\) positive, define
\(fS=fS_1-fS_2\).  This is independent of the decomposition: from a second
decomposition \(S=S'_1-S'_2\) one gets
\(S_1+S'_2=S'_1+S_2\); equality of these positive currents gives equality
of every attached measure, hence
\(fS_1+fS'_2=fS'_1+fS_2\).  The local definitions agree on overlaps and
glue, and the same computation proves compatibility with \(\iota_*\).  We
make no assertion for an arbitrary order-zero current.
This proves the assertions about currents and their weak topology.

For completeness, the psh assertion is not an additional theorem about
nilpotents.  At stage zero of the CLD filtration, smoothness and positivity
of \(\mathrm d'\mathrm d''u\) agree by
\eqref{eq:p21p3-nr-form-isomorphism}.  Successor stages take decreasing
limits, the continuous filtration takes local uniform limits, and limit
ordinal stages take local unions.  The underlying open sets and functions
are identical.  Transfinite induction in the filtration of
\cref{def:ch16-psh-filtration} therefore identifies the psh and continuous
psh sheaves at every stage, hence after stabilization.  This proves
\cref{itm:p21p3-nr:calculus}; the \(G\)-form argument is the same, using the
\(G\)-form variant stated immediately after
\cite[Proposition~15.1.3]{CLD25}.

\smallskip
\emph{The integration current and generic multiplicity.}
By \cite[Definition~13.2.5 and \S13.5.2]{CLD25}, the coefficient of
\([X_i]\) in the fundamental cycle is the positive generic multiplicity
\(m_i\), and the fundamental cycle is

\[
                         [X]=\sum_i m_i[X_i].
\]

By \cref{thm:current-fundamental-cycle}, which is
\cite[Proposition~16.6.3]{CLD25}, integration agrees with the current of
the fundamental cycle.  Applying the reduced identity
\eqref{eq:p21p3-fundamental-cycle} to \(X_{\mathrm{red}}\) gives

\[
 \delta_X=\sum_i m_i\delta_i,
 \qquad
 \iota_*\delta_{X_{\mathrm{red}}}=\sum_i\delta_i.
\]

This is \eqref{eq:p21p3-nr-fundamental-cycle}.  The components are locally
finite and all \(m_i\) are positive integers, so the sums are locally
finite.  Their equality forces \(m_i=1\) for every \(i\), by restriction to
an open set of \(X_i\) meeting no other component; the converse is
immediate.  This proves \cref{itm:p21p3-nr:integration}.

\smallskip
\emph{Classification.}
Let \(T\) be positive and closed on \(X\).  By
\eqref{eq:p21p3-nr-current-isomorphism} there is a unique positive closed
current \(T_{\mathrm{red}}\) on \(X_{\mathrm{red}}\) with
\(T=\iota_*T_{\mathrm{red}}\).  Apply
\cref{prop:p21p3-curve-rigidity}\,
\cref{itm:p21p3-curve:classification} to obtain

\[
                  T_{\mathrm{red}}=\sum_i c_i
                    (\widetilde X_i\to X_{\mathrm{red}})_*
                    \delta_{\widetilde X_i},
                  \qquad c_i\geq0.
\]

Pushing forward by \(\iota\) gives
\eqref{eq:p21p3-nr-classification}, with the same local finiteness and
uniqueness.  If \(X_{\mathrm{red}}\) is connected and rig-smooth there is
one component, say of generic multiplicity \(m\).  Then
\(\delta_X=m\delta_1\), and \(c_1\delta_1=(c_1/m)\delta_X\), proving the
last assertion of \cref{itm:p21p3-nr:classification}.

\smallskip
\emph{The multiplicity-weighted curvature identity.}
Fix \(W,U,V\) as in \cref{itm:p21p3-nr:bedford-taylor}, identify \(W\) with
the corresponding open subset of \(X_{\mathrm{red}}\), and put
\(\widetilde W_i=\nu_i^{-1}(W)\).  Denote the two finite base-changed maps by
\[
 \nu_{i,W}^{X}:\widetilde W_i\longrightarrow W\subset X,
 \qquad
 \nu_{i,W}^{\mathrm{red}}:\widetilde W_i\longrightarrow
                         W\subset X_{\mathrm{red}}.
\]
Define positive Radon measures

\[
 \mu_i^U:=\mathrm d'\mathrm d''[\nu_i^*U],
 \qquad
 \mu_i^V:=\mathrm d'\mathrm d''[\nu_i^*V],
 \qquad
 \sigma_i:=\mu_i^U-\mu_i^V.
\]

Pullback preserves continuous psh functions by
\cite[Proposition~8.1.5]{CLD25}, and positivity of the displayed measures
is \cite[Lemma~8.5.5]{CLD25}.  Restrict
\eqref{eq:p21p3-nr-fundamental-cycle} to \(W\).  For a continuous function
\(F\) on \(W\), use the definitions
\[
 [F]_X=F(\delta_X|_W),\qquad
 [F]_{X_{\mathrm{red}}}=F(\delta_{X_{\mathrm{red}}}|_W)
\]
from the statement.  Evaluating on a compactly supported test form and using
the defining transpose pairing for finite pushforward gives

\begin{align}
 [F]_X
   &=\sum_i m_i(\nu_{i,W}^{X})_*
       \bigl(((\nu_{i,W}^{X})^*F)\delta_{\widetilde W_i}\bigr),
 \label{eq:p21p3-nr-function-current-X}\\
 [F]_{X_{\mathrm{red}}}
   &=\sum_i(\nu_{i,W}^{\mathrm{red}})_*
       \bigl(((\nu_{i,W}^{\mathrm{red}})^*F)
                    \delta_{\widetilde W_i}\bigr).
 \label{eq:p21p3-nr-function-current-red}
\end{align}

For \(F=U,V\), the continuous-psh projection formula is precisely
\cite[Proposition~8.5.6]{CLD25}; equivalently, one may differentiate the
displayed identities and use commutation of proper pushforward with the
differentials.  Thus equality of the two curvature currents on \(X\) is
equivalent to

\begin{equation}
\label{eq:p21p3-nr-weighted-measure-sum}
                       \sum_i m_i(\nu_{i,W}^{X})_*\sigma_i=0,
\end{equation}

whereas equality on \(X_{\mathrm{red}}\) is equivalent to
\[
                       \sum_i
                         (\nu_{i,W}^{\mathrm{red}})_*\sigma_i=0.
\]

Let \(Z\subset X_{\mathrm{red}}\) be the union of the conductor locus of
the normalization and the pairwise intersections of distinct components,
and put
\[
 Z_W:=Z\cap W,
 \qquad
 Z_i:=\nu_i^{-1}(Z),
 \qquad
 Z_{i,W}:=\nu_i^{-1}(Z_W)=Z_i\cap\widetilde W_i.
\]
The proof of
\cref{prop:p21p3-curve-rigidity} shows that \(Z\) and all \(Z_i\) are
locally finite sets of type-\(1\) points, and that away from \(Z\) the
component images are disjoint and every \(\nu_i\) is an analytic
isomorphism.  Since \(m_i>0\), either the weighted or the unweighted
measure identity therefore implies

\[
                       \operatorname{Supp}\sigma_i\subseteq Z_{i,W}
                       \quad\text{for every }i.
\]

By \cref{lem:p21p3-type-one-zero-atom}, neither positive measure
\(\mu_i^U\) nor \(\mu_i^V\) has an atom at a point of \(Z_{i,W}\).  A signed
Radon measure supported on a locally finite set is the locally finite sum
of its point masses.  Hence

\begin{equation}
\label{eq:p21p3-nr-component-measure-equality}
                         \sigma_i=0\qquad\text{for every }i.
\end{equation}

Conversely, \eqref{eq:p21p3-nr-component-measure-equality} implies both
the weighted and the unweighted measure identities.  This proves the exact
curvature equivalence \eqref{eq:p21p3-nr-curvature-equivalence}; in
particular, the generic multiplicities in
\eqref{eq:p21p3-nr-weighted-measure-sum} cannot conceal cancellation at a
component intersection.

Finally let \(h=U-V\) and write \(T=\sum_i c_i\delta_i\).  Multiplication by
\(h\) is defined for the positive current \(T\) by
\cite[Example~7.2.12]{CLD25}; continuity of \(h\) makes it locally bounded.
Using proper pushforward and
\eqref{eq:p21p3-nr-component-measure-equality}, we obtain

\begin{align*}
 \mathrm d'\mathrm d''(hT|_W)
 &=\sum_i c_i(\nu_{i,W}^{X})_*
     \mathrm d'\mathrm d''
       [\nu_i^*U-\nu_i^*V]\\
 &=\sum_i c_i(\nu_{i,W}^{X})_*(\mu_i^U-\mu_i^V)=0.
\end{align*}

This is \eqref{eq:p21p3-nr-bt-conclusion} and is precisely the
Bedford--Taylor condition of \cite[Definition~8.6.2]{CLD25}.  Since the
same argument with all \(m_i=1\) is the reduced condition, the predicate is
preserved in both directions by \eqref{eq:p21p3-nr-current-isomorphism}.
This proves \cref{itm:p21p3-nr:bedford-taylor} and completes the proof.
\end{proof}
